% Section 2.5, from lectures/L02/appendix/b_exercises.md.

\section{Exercises}
\label{c2:sec:exercises}\label{c2:app:b}

\begin{aside}
\textbf{How to work these.} The Reflection exercise is answered on paper, from the chapter.
Exercise~\ref{c2:ex:1} needs no board: as its tip shows, a small program on your own computer checks
the expressions. The other Code exercises are programs for the board, and this is the chapter that
sets up the toolchain for them: build and flash them the way \secref{c2:app:c} (Microchip Studio)
or \secref{c2:app:d} (Linux) describes, or run them in the simulator of Appendix~\ref{sim:ch} if you
have no board yet.

\textbf{How to check your work.} The Reflection exercise, and the register values the expressions
of Exercise~\ref{c2:ex:1} should produce, are answered in Appendix~\ref{app:answers}. The Code
exercises have no printed solution: the board tells you whether a program does what the exercise
describes. Try each exercise before you read its answer.
\end{aside}

% ------------------------------------------------------------------------------------------
\exerciseset{Bit manipulation}
\label{c2:set:1}

\codeexercise{Set, clear and read a bit}
\label{c2:ex:1}
Given a register currently holding \code{0b10110010}, write single-line C expressions (using
\code{|}, \code{&}, \code{~}, \code{<<}) to:
\begin{parts}
  \item Set bit 0 without affecting any other bit.
  \item Clear bit 5 without affecting any other bit.
  \item Read whether bit 4 is set, as a boolean expression.
\end{parts}

\note[Tip] You don't need real hardware to test this; write a small program on your own machine and
check the results by invoking \code{printf()} from \code{<stdio.h>}:

\begin{ccode}
#include <stdint.h>
#include <stdio.h>

int main(void)
{
    const uint8_t reg = (1U << 3U);
    printf("%u\n", reg);
    return 0;
}
\end{ccode}

\codeexercise{LED macros}
\label{c2:ex:2}
A macro is defined as \code{#define LED 6U}. Write the two macros \code{LED_ON} and
\code{LED_OFF}, operating on a register called \code{PORTD}, following the OR/AND-NOT pattern used
in the chapter.

% ------------------------------------------------------------------------------------------
\exerciseset{Registers}
\label{c2:set:2}

\reflexercise{Which register?}
\label{c2:ex:3}
For each of the following, name the correct register (\code{DDRx}, \code{PORTx}, or \code{PINx})
and explain in one sentence why:
\begin{parts}
  \item Configuring D5 as an output.
  \item Reading whether a button on D2 is currently pressed.
  \item Enabling the internal pull-up resistor on D7 configured as an input.
\end{parts}

% ------------------------------------------------------------------------------------------
\exerciseset{Program}
\label{c2:set:3}

\codeexercise{Two buttons, two LEDs}
\label{c2:ex:4}
Wire up LEDs and buttons as described below:
\begin{itemize}
  \item \code{LED1} on D6 and \code{LED2} on D7.
  \item \code{BTN1} on D2 and \code{BTN2} on D3.
\end{itemize}
Write a complete program, from scratch, that:
\begin{itemize}
  \item Lights \code{LED1} if \code{BTN1} is pressed and turns it off otherwise.
  \item Lights \code{LED2} if \code{BTN2} is released and turns it off otherwise.
\end{itemize}
Don't copy the chapter's example directly; use these different pin numbers and re-derive the
macros yourself.

\codeexercise{Cycling three LEDs}
\label{c2:ex:5}
Wire up LEDs and buttons as described below:
\begin{itemize}
  \item \code{LED1} on D5, \code{LED2} on D6, and \code{LED3} on D7.
  \item \code{BTN1} on D4.
\end{itemize}
Write a complete program, from scratch, that:
\begin{itemize}
  \item While \code{BTN1} is held, cycles through the LEDs one at a time, in order (\code{LED1}
    $\rightarrow$ \code{LED2} $\rightarrow$ \code{LED3} $\rightarrow$ \code{LED1} \dots), lighting
    exactly one for 100\,ms before moving to the next.
  \item Turns off all LEDs otherwise.
\end{itemize}
